\section{Conclusion and Future Scope}\label{sec:conclusion}

\subsection{Conclusion}

This study asked whether, under customer-level SGCC theft labels and no verified theft-event timestamps, learning representations of unlabeled consumption windows before distribution-preserving aggregation improves customer-level discrimination beyond applying the same aggregation directly to handcrafted window features. The completed evidence answers this question affirmatively within the governed SGCC protocol, while retaining important limits on interpretation.

The experimental progression clarified why the final formulation was needed. Normal-only FC, CNN, and LSTM reconstruction errors provided weak customer discrimination under the evaluated anomaly formulations. Customer-supervised mean, maximum, and gated-attention MIL better aligned training with \texttt{FLAG}, but each compressed a customer bag into one latent vector and remained below the engineered benchmark. AMP-LDP instead preserved the mean, stabilised population standard deviation, 10th percentile, median, 90th percentile, and maximum of every learned latent coordinate. The matched DIRECT control applied those same statistics to handcrafted \textsc{summary\_17} features, while the original Stage~17B \textsc{dev\_train}-fitted FULL\_48 HGB supplied a strong engineered customer-profile comparator.

The frozen five-seed AMP-LDP ensemble replicated after development without refitting or seed selection. On the final evaluation cohort it achieved 0.811729 ROC-AUC and 0.421863 average precision. Relative to DIRECT, the paired improvements were 0.016762 ROC-AUC with a 95\% percentile-bootstrap interval of [0.006282, 0.027865] and 0.053424 average precision with an interval of [0.027428, 0.077907]. Both intervals remained above zero, and the qualitative learned-encoding advantage had the same direction in development and independent confirmation. AMP-LDP also exceeded the frozen Study~2 FULL\_48 comparator in the final paired analysis, while four of five individual seeds satisfied its preregistered parity margins and no seed-collapse condition fired.

These results support the bounded conclusion that learning window representations before distribution-preserving customer aggregation adds reproducible customer-level discrimination beyond identical direct aggregation across the governed development, confirmation, and final evaluation roles. They do not prove that a single architectural component causally produced the gain, that the effect size is fixed across cohorts, or that a high-scoring window is a verified theft period. Study~1 and Study~2 remain closed, and the report makes no claim of external generalisation, universal novelty, state-of-the-art superiority, or deployment readiness.

\subsection{Future Scope}

Future work should begin with external validation on independently collected smart-meter data under a prospectively frozen protocol. Such evaluation should preserve customer-level split integrity, report natural prevalence, and document preprocessing and label provenance so that transfer across utilities, regions, meter populations, and time periods can be assessed without relying on dataset-name equivalence.

If utilities can provide verified theft-event timestamps or inspected event intervals, a separate study could evaluate temporal localisation and determine whether latent windows correspond to confirmed events. Those labels would enable event-level sensitivity and false-alarm analysis; they should not be approximated by assigning a customer flag to every window. Independently preregistered sensitivity analyses could also examine alternative window lengths, overlap policies, summary features, aggregation statistics, encoder capacities, and chronological extensions. Their purpose would be to measure design dependence, not to retroactively rescue or revise the frozen AMP-LDP result, and no alternative should be presumed superior in advance.

Operational research should evaluate probability calibration and decision thresholds using utility-defined inspection costs, capacity constraints, and acceptable false-positive burdens. Prospective studies could report precision among the highest-ranked customers, inspections required per confirmed case, decision-curve or cost-sensitive outcomes, and stability under temporal or regional distribution shift. Finally, deployment-oriented profiling should measure inference latency, throughput, memory and energy use, monitoring requirements, and the maintenance implications of the five-seed ensemble relative to smaller alternatives. These extensions would determine practical suitability; the current study neither completes them nor guarantees favourable outcomes.
